\documentclass{article}
\usepackage[utf8]{inputenc}
\usepackage{graphicx} % Required for inserting images
\usepackage{amsmath}
\usepackage{booktabs}
\usepackage{float} % [H] placement: put a float exactly where it is in the source

\title{Title: TBD}
\author{}

\begin{document}

\maketitle

\section{Introduction}

\section{Results}

\subsection{Chord execution improved with practice}

Sixteen participants (XX females) practised multi-finger hand configurations combining flexion and extension around the metacarpophalangeal joint ("chords") over 24 daily sessions. Four chords were practised in each session (trained chords). The remaining 4 chords (untrained chords) were performed only in sessions 1 and 2 (pre-training sessions), sessions 3, 9 and 23 (MRI sessions) and sessions 4, 10 and 24 (behavioural testing sessions). The assignment of chords to the trained and untrained set was counterbalanced across participants, so that each chord was in the trained set for half of the participants. Each participant was assigned a unique set of trained and untrained chords (i.e., 16 different trained-untrained chord sets we used in the experiment out of $\binom{8}{4}=70$ possible combinations).\\

Initially, chord execution was challenging and participants failed to achieve the required configuration in most trials; by session 3 (i.e., the first MRI session), success rate raised >60\% (FIG) with most unsuccessful trials representing near-misses (FIG). At the end of the training period, both trained and untrained chords could be achieved with nearly 100\% success rate. This confirms our previous finding that all chords can be achieved without hard biomechanical limitations using finger forces from 2-5 N (REF). Gaining familiarity with the task through practice determined a general improvement in the ability to achieve chords, as shown by the similar success rate for trained and untrained chords in sessions 9-10 and 23-24. \\

Yet, chord execution was different between the trained and untrained set after practice. Trained chords were executed faster and with more synchronous finger presses, as showed by significantly lower execution time (ET) and mean deviation (MD) in sessions 9-10 and 23-24. This indicates that, while a general familiarity with the task is sufficient to achieve all chords, practice produces chord-specific performance gains. \\

\subsection{The execution of trained chords requires less activation in cortical motor regions}

On sessions 3, 9 and 23, participants performed trained and untrained chords while being scanned with 7T fMRI. The inspection of BOLD activity time-aligned with chord execution and averaged across participants within each region-of-interest (ROI; FIG) showed that our task elicted a strong activation across cortical motor areas, and suggested smaller responses in trained vs. untrained chords in sessions 9 and 23. To characterise BOLD activation in more detail, we fitted a general linear model (GLM) with separate regressors for each chord and MRI session to the BOLD time series recorded in each voxel. In all ROIs, the mean BOLD activation decreased overall across sessions, with a more pronounced decrease for trained vs. untrained chords (FIG). Confirming this observation, a repeated-measures ANOVA showed a significant main effect of SESSION ($F_{2,30} \geq 4.67$, $p \leq .024$) and CHORD SET ($F_{1,15} \geq 12.91$, $p \leq .003$), and a significant SESSION $\times$ CHORD SET interaction ($F_{2,30} \geq 7.11$, $p \leq .011$) in all ROIs (Table~\ref{tab:activation_stats}). In sessions 9 and 23, activation was lower for trained than for untrained chords. This difference was significant in all ROIs except M1 in session 9 (paired $t$-tests, $t_{15} \leq -2.14$, $p \leq .0495$; M1: $t_{15} = -1.88$, $p = .080$), and in all ROI in session 23 ($t_{15} \leq -3.43$, $p \leq .004$), with the strongest effects in PMd, SPLa, and SPLp. In contrast, activation did not differ between trained and untrained chords in session 3 in any ROI ($|t_{15}| \leq 1.08$, $p \geq .30$).

\begin{table}[H]
    \centering
    \resizebox{\textwidth}{!}{% scale the table to the text width
    \begin{tabular}{lccccccc}
        \toprule
         & SMA & PMd & PMv & M1 & S1 & SPLa & SPLp \\
        \midrule
        \multicolumn{8}{l}{\textit{Repeated-measures ANOVA}} \\
        SESSION, $F_{2,30}$ & 8.51 (.002) & 9.48 (.002) & 13.82 ($<$.001) & 4.67 (.024) & 6.12 (.014) & 12.69 ($<$.001) & 11.74 (.001) \\
        CHORD SET, $F_{1,15}$ & 13.45 (.002) & 26.08 ($<$.001) & 14.30 (.002) & 19.89 ($<$.001) & 12.91 (.003) & 22.11 ($<$.001) & 16.66 ($<$.001) \\
        SESSION $\times$ CHORD SET, $F_{2,30}$ & 7.41 (.011) & 15.54 ($<$.001) & 8.89 (.005) & 7.11 (.005) & 12.90 ($<$.001) & 19.69 ($<$.001) & 19.57 ($<$.001) \\
        \midrule
        \multicolumn{8}{l}{\textit{Paired $t$-tests, trained vs. untrained}} \\
        Session 3, $t_{15}$ & 1.08 (.299) & 0.31 (.763) & $-$0.58 (.573) & $-$0.69 (.499) & $-$0.42 (.681) & 0.23 (.821) & 0.65 (.528) \\
        Session 9, $t_{15}$ & $-$2.33 (.034) & $-$4.17 ($<$.001) & $-$2.90 (.011) & $-$1.88 (.080) & $-$2.14 (.0495) & $-$3.86 (.002) & $-$3.53 (.003) \\
        Session 23, $t_{15}$ & $-$3.43 (.004) & $-$5.30 ($<$.001) & $-$3.65 (.002) & $-$4.39 ($<$.001) & $-$4.57 ($<$.001) & $-$5.10 ($<$.001) & $-$4.90 ($<$.001) \\
        \bottomrule
    \end{tabular}}
    \caption{ROI-based statistics of the univariate activation in the left hemisphere. Uncorrected $p$-values are provided.}
    
    \label{tab:activation_stats}
\end{table}

\subsection{Chord information across cortical motor regions remains stable with practice}

Univariate activation offers a partial view of the information encoded in BOLD activity and on the changes that occurred after practice. We used representational similarity analysis (RSA) to assess multi-voxel activity patterns and characterise how the information about the task is encoded across cortical motor regions. First, we evaluated if BOLD activity patterns differed significantly between chords. To this end, we calculated the cross-validated Mahalanobis (crossnobis) dissimilarities between the BOLD activity patterns elicited during the execution of each chord. In a region without chord representation, the average dissimilarity estimate should be zero. In contrast, all ROI showed above-chance encoding already in session 3. \\

Next, we sought to establish whether the information represented in each ROI remained stable across participants. To this end, we correlated the pairwise crossnobis dissimilarities of each participant with the average dissimilarities of the remaining participants (lower noise ceiling) and of the entire group (upper noise ceiling). The lower noise ceiling reflects the part of each participant's geometry that is shared with the rest of the group, whereas the upper noise ceiling reflects the maximum correlation attainable given the measurement noise in each participant's data. Already in session 3, the lower noise ceiling was significantly greater than zero in all ROIs (all t(15) $\geq 2.7$, p < .01). In most ROIs it was also close to the upper noise ceiling, reaching 75–95\% of its value. This suggests that the information that cortical motor regions encode about chords is consistent across participants, and that individual differences are small, in line with previous work showing that the representational geometry of finger movements is stable across individuals (Ejaz et al., 2015). \\

What chord information is captured by brain activity? To investigate this, we defined eight candidate representation models that predicted the geometry of BOLD activity patterns for each chord in voxel space based on different chord properties. (i) The \textit{identity} model predicted that BOLD patterns for each chords were equally distance in voxel space. (ii) The \textit{single finger} model predicted that the distance between BOLD activity patterns for each chord reflected the distance between the patterns obtained by summing the activity patterns evoked by the flexion or extension of each individual finger involved in the chord, measured in the same ROI in a separate dataset (Arbuckle et al., 2020). (iii) The \textit{muscle pattern} model predicted that the distance between BOLD activity patterns reflected the distance between the corresponding muscle activity patterns recorded with EMG outside of the MRI scanner in a separate group of participants (see Methods). (iv) The \textit{hand use} model predicted that chords requiring finger movements that frequently co-occur in everyday hand use evoke similar BOLD activity patterns, based on kinematic recordings of everyday hand activities (Ingram et al., 2008). (v) The \textit{signed force pattern} model predicted that the distance between BOLD activity patterns reflected the distance between the patterns of signed forces exerted by the five fingers while executing each chord, averaged across the other participants. Three further models were based on behavioural measures obtained from the other participants on session 4 (i.e., the last session before training started). (vi) The \textit{mean deviation} model predicted that the distance between BOLD activity patterns reflected differences in chord difficulty, as indexed by the mean deviation of the finger forces from the ideal straight force trajectory. (vii) The \textit{flexion-extension bias} model predicted that distance between BOLD activity patterns reflected whether a chord leaned towards flexion or extension, as indexed by the mean signed force. (viii) The \textit{mean force magnitude} model predicted that distance between BOLD activity patterns reflected the overall force required to produce the chord. We then used Pattern Component Modeling (PCM) to evaluate the likelihood of the observed BOLD activity patterns in each ROI under the hypothesis that their representational geometry corresponds to a weighted combination of these models. \\

Consistent with above-chance encoding (as indexed by crossnobis dissimilarities > 0), in session 3 the likelihood of the BOLD activity patterns under the identity model was significantly higher than under a null model that predicted no difference between activity patterns in all ROIs (all t\textit{s}(15) $\geq$ 4.42, all p\textit{s} < .001). We then asked whether the activity patterns also contained more structured information about chords beyond the fact that each chord evoked a distinct pattern. To this end, we fitted weighted combinations of each other candidate models plus the identity model compared the likelihood of the data under each of these combined model vs. under the identity alone, correcting for the additional parameter (AIC). All combination models outperformed identity alone in all ROIs (all t\textit{s}(15) $\geq$ 2.12, all p\textit{s} < .03; TAB), with the exception of the model containing \textit{mean deviation}, which outperformed identity only in S1, SPLa and SPLp. 

Importantly, the candidate models are not independent of each other but predict partly overlapping representational geometries (Fig. X). A model can therefore outperform identity alone simply because it shares a similar structure with another model. To isolate the unique contribution of each model, we fitted all combinations of the candidate models and, for each model, computed the log Bayes factor between the combinations that included it and those that did not, averaged across all combinations of the other models (see Methods). Because every combination included the identity model, this quantifies the evidence that a model explains part of the similarity structure that none of the other models accounts for. 

[RESULTS: e.g., The \textit{single finger} model made a unique contribution in PMd and S1 in all sessions, and in PMv and M1 after practice, whereas … ] This shows that, although several chord properties predict the geometry of the activity patterns, only a subset contributes uniquely, with the summed single-finger patterns accounting for most of the shared structure.


Finally, the \textit{trained} and \textit{untrained} models predicted that the BOLD activity patterns of the trained or untrained chords, respectively, were more distinct from each other and from the remaining chords than predicted by the other models.

... and then used PCM to evaluate the likelihood that the observed BOLD activity patterns come from ...






\subsection{Learning-related changes in BOLD activity are not the result of changes in the force output}
BOLD activity is influenced by the amount of force generated by the fingers (REF) and, more importantly, by the rate of force variation over time (REF). Therefore, learning-related changes in BOLD activity could reflect a genuine modification of chord encoding across cortical motor regions or be the product of how these two aspects of chord execution are altered after extensive practice. The average absolute force across fingers increased slightly (within the instructed 2-5 N force range) with practice, and roughly equally for trained and untrained chords (Fig.~\ref{fig:force_measures}a). On the other hand, the absolute force derivative showed an overall decrease with practice, which was more pronounced for trained vs. untrained chords in sessions 9 and 23 (Fig.~\ref{fig:force_measures}b), compatible with a more stable force output and fewer corrective force adjustments needed to maintain the force pattern required by the instructed chord.\\

\begin{figure}
    \centering
    \includegraphics[scale=.5]{figures/force_measures.pdf}
    \caption{Enter Caption}
    \label{fig:force_measures}
\end{figure}

The changes in the absolute force derivative averaged across fingers showed a different time course the univariate activation across cortical motor areas (FIG). In the absolute force derivative, the difference between trained and untrained chords was roughly similar in sessions 9 and 23 (Fig.~\ref{fig:force_measures}b), while the difference in univariate activation kept increasing in most ROI (FIG). Therefore, it is unlikely that the decrease in univariate activation with practice is simply the result of the decrease in the absolute force derivative. To demonstrate this quantitatively, we performed a linear regression analysis between the difference in univariate activation for trained and untrained chords and the corresponding average difference in absolute force derivative in each participant. If the difference in univariate activation is (at least partly) independent of the difference in force derivative between trained and untrained chords, the regression intercept should be significantly $<0$, meaning that a hypothetical participant with no difference in absolute force derivative would still show reduced activation for trained vs. untrained chords. [RESULTS]\\

In addition to univariate activation, changes in the force output after practice could also influence multi-voxel activity patterns across cortical motor areas and potentially drive the pattern separation that we observed in M1 and S1. In fact, the reduction in the absolute force derivative in sessions 9 and 23 was not only larger in the trained vs. untrained chords (Fig.~\ref{fig:force_measures}b), but in the trained chords was also larger in the finger that had to remain neutral relative to the fingers that were activated in flexion or extension (Fig.~\ref{fig:force_measures}c). This indicates that with practice participants learned to perform trained chords with fewer force adjustments on the uninstructed finger. Because the uninstructed finger typically differed across chords, this selective reduction led to an increase in the angular distance between absolute force derivative patterns for trained vs. untrained chords in sessions 9 and 23 (Fig.~\ref{fig:force_measures}d). Therefore, the increased separation between neural activity patterns for trained vs. untrained chords in M1 and S1 could, in principle, simply reflect the different way in which trained chords were executed after practice. \\

The angular distance between absolute force derivative patterns for trained chords showed a different time course compared to the angular distance between corresponding neural activity patterns (FIG), making it unlikely that this change in the motor output could be alone responsible for the separation between neural patterns that occurred with practice. We performed three separate control analyses to quantitatively support this conclusion. \\

First, we fitted a linear model between pairwise cosine dissimilarities in absolute force derivative patterns and BOLD dissimilarities in session 3. Then, we used the fitted coefficients to predict the BOLD cosine dissimilarities in each session from the corresponding absolute force derivative dissimilarities, and transformed the residual dissimilarities not captured by the model into angular distances. If BOLD dissimilarities in sessions 9 and 23 are simply driven by force derivative dissimilarities, then the residual angular distances not captured by the model in sessions 9 and 23 should not be significantly $>0$. Instead, at least in session 23, the residual angular distance between trained chords in M1, S1 and SPLa was significantly larger than between untrained chords (Fig.~\ref{fig:control_ana_pattern}a).\\

\begin{figure}
    \centering
    \includegraphics[scale=.5]{figures/control_ana_pattern.pdf}
    \caption{Enter Caption}
    \label{fig:control_ana_pattern}
\end{figure}

Second, we modelled the cosine BOLD dissimilarities in each session as a function of the corresponding absolute force derivative dissimilarities. The absolute force derivative coefficient was positive in M1 and S1 in sessions 9 and 23, confirming that the force derivative explained part of the BOLD geometry (Fig.~\ref{fig:control_ana_pattern}b). We then removed the contribution of the absolute force derivative from the cosine dissimilarities and converted the force-adjusted cosine dissimilarities into angular distances. After adjusting for absolute force derivative, the angular distance in session 23 remained larger for trained than untrained chords in M1, S1 and SPLa (Fig.~\ref{fig:control_ana_pattern}c).\\

Third, in each session and ROI, we regressed the difference in angular distance between neural activity patterns for trained vs. untrained chords onto the corresponding difference in angular distance between absolute force derivative patterns. Similar to the regression analysis that we used for univariate activation, if the separation of neural activity patterns between trained chords is (at least partly) independent of the separation between absolute force derivative patterns, the regression intercept should be significantly $>0$, meaning that a hypothetical participant with no difference in absolute force derivative patterns for trained and untrained chords would still show a larger angular distance between neural patterns for trained chords. In session 23, the intercept was generally positive across cortical motor regions, reaching significance in M1 (Fig.~\ref{fig:control_ana_pattern}d).

\section{Methods}

\subsection{Participants}

\subsection{Task}

\subsection{Experimental procedures}

\subsection{Evaluation of behavioural confounds on pattern separation}

We used three separate approaches to evaluate whether the increased separation between neural activity patterns for trained vs. untrained chords across cortical motor regions was simply the result of changes in motor performance after practice. \\

First, we modelled the cosine BOLD dissimilarities on session 3 for each participant and ROI as a function of the corresponding dissimilarities between absolute force derivative patterns across fingers, as follows:

\begin{equation}
    \mathbf{y}_3 = \mathbf{X}_3 \boldsymbol{\beta} + \boldsymbol{\varepsilon}_3,
    \qquad
    \mathbf{X}_t = \begin{bmatrix} \mathbf{1} & \mathbf{f}_3 & \bar{\mathbf{y}}_3 \end{bmatrix} \in \mathbb{R}^{12 \times 3}
%\label{eq:force_prediction}
\end{equation}

where $\mathbf{y}_3 \in \mathbb{R}^{12}$ is the 12-dimensional vector of pairwise cosine dissimilarities between BOLD activity patterns for trained and untrained chords (6 trained and 6 untrained pairs) on session 3; $\mathbf{f}_3 \in \mathbb{R}^{12}$ is the corresponding vector of pairwise cosine dissimilarities between absolute force derivative patterns; and $\bar{\mathbf{y}}_3$ is the vector of group-average pairwise cosine dissimilarities between BOLD activity patterns in session 3. The group term captures the pair-specific geometry of each region that is unrelated to force, and is identical across sessions. The fitted coefficients were then used to predict the BOLD dissimilarities in each session, and the residuals were computed as:
\begin{equation}
    \hat{\mathbf{y}}_t = \mathbf{X}_t \hat{\boldsymbol{\beta}},
    \qquad
    \mathbf{r}_t = \mathbf{y}_t - \hat{\mathbf{y}}_t,
    \qquad t \in \{3, 9, 23\}.
\end{equation}
Since $\bar{\mathbf{y}}_3$ does not change across sessions, the model can only predict a change in BOLD dissimilarities through a change in $\mathbf{f}_t$. The residuals $\mathbf{r}_t$ therefore reflect the change in neural geometry not accounted for by the change in absolute force derivates. Residuals were averaged across the pairs of trained and of untrained chords for each participant, and compared between chord sets using [TEST]. \\

Second, we modelled the cosine BOLD dissimilarities in each session as a function of the group-average geometry in session 3, the absolute force derivative dissimilarities and the chord set (i.e., trained vs. untrained), as follows:
\begin{equation}
    \begin{bmatrix} \mathbf{y}_3 \\ \mathbf{y}_9 \\ \mathbf{y}_{23} \end{bmatrix}
    =
    \underbrace{\begin{bmatrix}
        \bar{\mathbf{y}}_3 & \mathbf{Z}_3 & \mathbf{0} & \mathbf{0} \\
        \bar{\mathbf{y}}_3 & \mathbf{0} & \mathbf{Z}_9 & \mathbf{0} \\
        \bar{\mathbf{y}}_3 & \mathbf{0} & \mathbf{0} & \mathbf{Z}_{23}
    \end{bmatrix}}_{\mathbf{X} \,\in\, \mathbb{R}^{36 \times 10}}
    \begin{bmatrix} \beta_G \\ \boldsymbol{\theta}_3 \\ \boldsymbol{\theta}_9 \\ \boldsymbol{\theta}_{23} \end{bmatrix}
    + \boldsymbol{\varepsilon},
    \label{eq:ancova}
\end{equation}

where $\mathbf{y}_t \in \mathbb{R}^{12}$ is the 12-dimensional vector of pairwise cosine dissimilarities between BOLD activity patterns for trained and untrained chords (6 trained and 6 untrained pairs) on session $t$; $\bar{\mathbf{y}}_3 \in \mathbb{R}^{12}$ is the vector of group-averaged cosine dissimilarities for the same pairs of BOLD activity patterns on session 3, with a single slope $\beta_G$ shared across sessions; and $\mathbf{Z}_t = \begin{bmatrix} \mathbf{f}_t & \mathbf{c} & \mathbf{1} \end{bmatrix} \in \mathbb{R}^{12 \times 3}$ contains the session-specific regressors, with coefficients $\boldsymbol{\theta}_t = [\gamma_t, \delta_t, \alpha_t]^\top$. Here, $\mathbf{f}_t$ is the vector of cosine dissimilarities between absolute force derivative patterns, $\mathbf{c} \in \{-1, +1\}^{P}$ codes chord set ($+1$ trained, $-1$ untrained). The model includes a session-specific intercept. We then computed the force-adjusted dissimilarities as:
\begin{equation}
    \tilde{\mathbf{y}}_t = \mathbf{y}_t - \hat{\beta}_G\left(\bar{\mathbf{y}}_3 - \langle \bar{\mathbf{y}}_3 \rangle\right) - \hat{\gamma}_t\left(\mathbf{f}_t - \langle \mathbf{f} \rangle\right),
\end{equation}
where $\langle\cdot\rangle$ denotes the mean over all pairs and sessions of the participant, so that adjusted values remain on the scale of the original dissimilarities. The force-adjusted cosine dissimilarities were converted into angular distances, $\tilde{\boldsymbol{\phi}}_t = \arccos(1 - \tilde{\mathbf{y}}_t)$, averaged across trained and untrained pairs in each session, and compared with one-sided paired $t$-tests across participants (trained $>$ untrained). The estimated force slopes $\hat{\gamma}_t$ were tested against zero with one-sample $t$-tests.

Third, for each session and participant, we computed the difference between the average angular distance in BOLD and absolute force derivative patterns for trained and untrained chords ($\Delta\phi^{\mathrm{BOLD}}_{n,t}$ and $\Delta\phi^{\mathrm{force}}_{n,t}$, respectively). Then, we fitted, separately for each session and ROI, a linear model that predicted the average difference in angular distance between BOLD activity patterns for trained and untrained chords $\Delta\phi^{\mathrm{BOLD}}_{n,t}$ as a function of the difference in angular distance between absolute force derivative $\Delta\phi^{\mathrm{force}}_{n,t}$, as follows:
\begin{equation}
    \Delta\phi^{\mathrm{BOLD}}_{n,t} = \beta_0 + \beta_1\,\Delta\phi^{\mathrm{force}}_{n,t} + \varepsilon_{n,t}
\end{equation}
The intercept $\beta_0$ estimates the difference in angular distance between BOLD activity patterns for trained and untrained chords for a participant with no difference in force derivative geometry, and was tested against zero with a one-sided $t$-test.

⁄\end{document}